% ৪.০ অধ্যায় সূচনা: ওয়েব ডিজাইন পরিচিতি ও HTML
\section{অধ্যায়ের সূচনা: ওয়েব ডিজাইন পরিচিতি ও HTML}

\begin{learningbox}
এই অধ্যায়ে শিক্ষার্থী ওয়েব ডিজাইনের মূল ধারণা, ওয়েবসাইটের কাঠামো এবং HTML ব্যবহার করে সহজ একটি ওয়েব পেজ কিভাবে তৈরি করতে হয় তা শিখবে। practical অংশে image, hyperlink, table ও form ব্যবহার করে একটি basic page তৈরি ও publish-এর ধারণা দেয়া হবে।
\end{learningbox}

\subsection{অধ্যায়ের নাম}
\begin{definitionbox}{অধ্যায়ের নাম}
চতুর্থ অধ্যায়: \textbf{ওয়েব ডিজাইন পরিচিতি ও HTML} (Introduction to Web Design and HTML)
\end{definitionbox}

\subsection{কেন এই অধ্যায়?}
বর্তমান যুগ তথ্য ও যোগাযোগ প্রযুক্তির যুগ। আমরা প্রতিদিন ইন্টারনেট ব্যবহার করে তথ্য সংগ্রহ করি, অনলাইন ফর্ম পূরণ করি এবং বিভিন্ন ওয়েবসাইট থেকে সেবা গ্রহণ করি। এই অধ্যায়ে শেখানো বিষয়গুলো দিয়ে শিক্ষার্থী সহজেই একটি বাস্তবসম্মত ও ব্যবহারবান্ধব ওয়েব পেজ তৈরি বা বুঝতে পারবে।

\begin{keypointbox}{Learning outcomes}
\begin{itemize}
\item ওয়েব ডিজাইনের ধারণা ও গুরুত্ব ব্যাখ্যা করতে পারবে।
\item ওয়েব পেজ, ওয়েবসাইট, হোম পেজ, সার্ভার ও ব্রাউজার-এর পার্থক্য বুঝবে।
\item HTML tag, element ও attribute চিনবে এবং সরল HTML document লিখতে পারবে।
\item image, hyperlink, list, table ও form ব্যবহার করে একটি simple page সাজাতে পারবে।
\end{itemize}
\end{keypointbox}

\subsection{Chapter flow}
\begin{center}
\fbox{\parbox{0.9\textwidth}{
Web Design Concept $\rightarrow$ Web Terms $\rightarrow$ Website Types $\rightarrow$ Website Structure $\rightarrow$ HTML Basic $\rightarrow$ Tags $\rightarrow$ Formatting $\rightarrow$ List $\rightarrow$ Hyperlink $\rightarrow$ Image $\rightarrow$ Table $\rightarrow$ Form $\rightarrow$ Mini Project $\rightarrow$ Publishing
}}
\end{center}

\subsection{Chapter keywords}
\begin{center}
\begin{tabular}{|p{0.28\textwidth}|p{0.60\textwidth}|}
\hline
Web Page & ওয়েব সার্ভারে রাখা একটি document/page \\
\hline
Website & একই domain-এর অধীনে একাধিক web page-এর সমষ্টি \\
\hline
Home Page & website খুললে প্রথম যে page দেখা যায় \\
\hline
WWW & World Wide Web \\
\hline
Browser & web page দেখার software \\
\hline
HTML & HyperText Markup Language \\
\hline
Tag/Element/Attribute & HTML structure-এর মূল উপাদান \\
\hline
Hyperlink & এক page থেকে অন্য page-এ যাওয়ার link \\
\hline
\end{tabular}
\end{center}

\subsection{Exam focus}
\begin{itemize}
\item HTML tag ও basic document structure
\item Hyperlink, Image, Table, Form-এর syntax ও ব্যবহার
\item Website structure ও static vs dynamic website
\end{itemize}

\begin{examplebox}{Real-life example}
তোমার কলেজ website-এ Notice, Result, Admission, Teacher Info, Contact page থাকবে; এগুলো মিলে একটি practical website তৈরি করে।
\end{examplebox}

\begin{practicebox}
\begin{enumerate}
\item লিখে দেখাও এই chapter-এ কী শেখা যাবে (৩ বিন্দু)।
\item ছোট একটি HTML skeleton লেখো এবং local browser-এ চালিয়ে দেখো।
\end{enumerate}
\end{practicebox}
